\chapter{Solution of Equations}

\section{Roots and Locating Roots}

\begin{definition}[Root or zero]
A number $r$ is a root of the equation $f(x)=0$ if
\[
 f(r)=0.
\]
The terms \emph{root}, \emph{solution} and \emph{zero of $f$} are used interchangeably.
\end{definition}

Linear and quadratic equations have explicit formulas.  For example,
\[
 ax+b=0\quad\Longrightarrow\quad x=-\frac ba,
 \qquad a\ne0,
\]
and
\[
 ax^2+bx+c=0
 \quad\Longrightarrow\quad
 x=\frac{-b\pm\sqrt{b^2-4ac}}{2a}.
\]
For general nonlinear equations, exact formulas are usually unavailable and numerical iteration is required.

\begin{example}
Find all roots of $x^3-4x=0$.
\end{example}

\begin{solution}[2.2cm]
Factor:
\[
 x^3-4x=x(x^2-4)=x(x-2)(x+2).
\]
Hence the roots are $-2$, $0$ and $2$.
\end{solution}

\begin{theorem}[Intermediate value theorem]
Let $f$ be continuous on $[a,b]$.  If $f(a)$ and $f(b)$ have opposite signs, then there exists at least one $r\in(a,b)$ such that $f(r)=0$.
\end{theorem}

The theorem guarantees existence but not uniqueness.  It is the theoretical basis of bracketing methods.

\begin{example}
Show that $x=\cos x$ has at least one solution in $(0,\pi/2)$.
\end{example}

\begin{solution}[2.8cm]
Let $f(x)=x-\cos x$.  Then $f$ is continuous and
\[
 f(0)=-1<0,
 \qquad
 f\left(\frac\pi2\right)=\frac\pi2>0.
\]
The intermediate value theorem therefore gives at least one root in $(0,\pi/2)$.  Trigonometric arguments must be measured in radians.
\end{solution}

\section{Bracketing Methods}

A bracketing method begins with two points $x_l$ and $x_u$ such that
\[
 f(x_l)f(x_u)<0.
\]
Continuity then guarantees a root inside the interval.  The interval is repeatedly reduced while preserving the sign change.

\subsection{Bisection Method}

\begin{algorithm}[Bisection]
\begin{enumerate}
  \item Choose $x_l$ and $x_u$ such that $f(x_l)f(x_u)<0$.
  \item Compute the midpoint
  \[
    x_r=\frac{x_l+x_u}{2}.
  \]
  \item If $f(x_l)f(x_r)<0$, set $x_u=x_r$.  If $f(x_l)f(x_r)>0$, set $x_l=x_r$.  If $f(x_r)=0$, stop.
  \item Repeat until the required tolerance is achieved.
\end{enumerate}
\end{algorithm}

The absolute interval error after $n$ bisections is bounded by
\[
 |r-x_r^{(n)}|\le\frac{b-a}{2^{n+1}}.
\]
A commonly used estimated percentage relative error is
\[
 \varepsilon_a
 =\left|\frac{x_r^{(n)}-x_r^{(n-1)}}{x_r^{(n)}}\right|100\%.
\]

\begin{example}
Use bisection to solve
\[
 x^{10}-1=0
\]
with initial interval $[0,1.3]$.  Stop when the estimated percentage error is below $8\%$.
\end{example}

\begin{solution}[7.0cm]
Let $f(x)=x^{10}-1$.  The bisection table is
\[
\begin{array}{c|ccc|c}
\toprule
n&x_l&x_u&x_r&\varepsilon_a(\%)\\
\midrule
1&0.0000&1.3000&0.6500&--\\
2&0.6500&1.3000&0.9750&33.3333\\
3&0.9750&1.3000&1.1375&14.2857\\
4&0.9750&1.1375&1.05625&7.6923\\
\bottomrule
\end{array}
\]
At the fourth iteration, $\varepsilon_a<8\%$.  Hence
\[
 r\approx1.0563.
\]
The exact positive root is $1$.
\end{solution}

Bisection is guaranteed to converge for a continuous function with a genuine sign change.  Its main disadvantages are slow linear convergence and the need for a valid bracket.  It may miss roots of even multiplicity because the function does not change sign across such roots.

\section{Open Methods}

Open methods do not maintain a bracketing interval.  They are generally faster when they converge, but a poor initial guess can lead to divergence or convergence to an unintended root.

\subsection{Fixed-Point Iteration}

To solve $f(x)=0$, rearrange the equation into
\[
 x=g(x)
\]
and iterate
\[
 x_{n+1}=g(x_n).
\]
A root of $f$ is then a fixed point of $g$.

\begin{example}
The function
\[
 f(x)=x^2-3x+e^x-2
\]
has a negative root.  Use fixed-point iteration with
\[
 x_{n+1}=\frac{x_n^2+e^{x_n}-2}{3},
 \qquad x_0=-0.5,
\]
to approximate the smaller root.
\end{example}

\begin{solution}[6.0cm]
The first few iterates are
\[
\begin{array}{c|r}
\toprule
n&x_n\\
\midrule
0&-0.5000000000\\
1&-0.3811564468\\
2&-0.3905495820\\
3&-0.3902620485\\
4&-0.3902720192\\
5&-0.3902716747\\
6&-0.3902716861\\
7&-0.3902716862\\
\bottomrule
\end{array}
\]
Thus the smaller root is
\[
 r\approx-0.3902716862.
\]
\end{solution}

Different rearrangements of the same equation can have very different convergence behaviour.  For example, a formula may be algebraically correct but have $|g'(r)|>1$, causing nearby errors to grow rather than shrink.

\begin{example}
Consider
\[
 x_{n+1}=\frac{3-x_n^3}{6}
\]
for solving $x^3+6x-3=0$.  Compare the behaviour for $x_0=-0.5$, $1.5$, $2.5$ and $5.5$.
\end{example}

\begin{solution}[6.0cm]
The first three starting values eventually approach
\[
 r\approx0.4814056.
\]
For instance,
\[
 -0.5,\ 0.52083,\ 0.47645,\ 0.48197,\ldots
\]
and
\[
 2.5,\ -2.10417,\ 2.05271,\ -0.94155,\ 0.63912,\ldots
\]
both settle toward the fixed point.  By contrast, $x_0=5.5$ produces
\[
 5.5,\ -27.2292,\ 3365.24,\ -6.35\times10^9,\ldots,
\]
so the iteration diverges.  Local convergence does not imply convergence from every starting value.
\end{solution}

\begin{theorem}[Convergence of fixed-point iteration]
Suppose $g$ is continuous on $[a,b]$ and maps $[a,b]$ into itself.  Then $g$ has at least one fixed point in $[a,b]$.  If, in addition, $g$ is differentiable and there is a constant $M<1$ such that
\[
 |g'(x)|\le M
 \qquad\text{for all }x\in[a,b],
\]
then the fixed point is unique and the iteration converges to it for every $x_0\in[a,b]$.
\end{theorem}

\subsection{Newton--Raphson Method}

Newton's method replaces the graph of $f$ by its tangent line at the current approximation.  The update is
\[
 x_{n+1}=x_n-\frac{f(x_n)}{f'(x_n)}.
\]
For a simple root and a sufficiently close starting value, convergence is usually quadratic.

\begin{algorithm}[Newton--Raphson]
\begin{enumerate}
  \item Choose an initial approximation $x_0$.
  \item Compute
  \[
    x_{n+1}=x_n-\frac{f(x_n)}{f'(x_n)}.
  \]
  \item Stop when $|x_{n+1}-x_n|<\varepsilon$, when $|f(x_{n+1})|$ is small, or when a maximum iteration count is reached.
\end{enumerate}
\end{algorithm}

\begin{example}
Use Newton's method to solve
\[
 x-e^{-x}=0
\]
for the root in $(0,2)$.  Continue until two successive approximations differ by less than $10^{-8}$.
\end{example}

\begin{solution}[5.8cm]
Take $x_0=1$.  Since
\[
 f(x)=x-e^{-x},
 \qquad
 f'(x)=1+e^{-x},
\]
we use
\[
 x_{n+1}=x_n-\frac{x_n-e^{-x_n}}{1+e^{-x_n}}.
\]
The iterates are
\[
\begin{array}{c|c|c}
\toprule
n&x_n&|x_n-x_{n-1}|\\
\midrule
1&0.5378828427&0.4621171573\\
2&0.5669869914&0.0291041487\\
3&0.5671432860&0.0001562946\\
4&0.5671432904&4.42\times10^{-9}\\
\bottomrule
\end{array}
\]
Thus
\[
 r\approx0.56714329.
\]
\end{solution}

Newton's method can fail if $f'(x_n)=0$, if the derivative is very small, or if the starting point lies in an unfavourable basin of attraction.

\subsection{Secant Method}

The secant method replaces the derivative by a finite-difference slope:
\[
 x_{n+1}
 =x_n-f(x_n)
 \frac{x_n-x_{n-1}}{f(x_n)-f(x_{n-1})}.
\]
It requires two initial values but no derivative.  Its convergence order is approximately
\[
 \frac{1+\sqrt5}{2}\approx1.618.
\]

\begin{example}
Use three secant iterations to approximate $\sqrt2$, with $f(x)=x^2-2$, $x_0=1.5$ and $x_1=1$.
\end{example}

\begin{solution}[3.5cm]
The successive approximations are
\[
 x_2=1.400000,
 \qquad
 x_3=1.4166667,
 \qquad
 x_4=1.4142012.
\]
Thus after three iterations,
\[
 \sqrt2\approx1.4142.
\]
\end{solution}

\section{Order of Convergence and Acceleration}

\begin{definition}[Order of convergence]
Suppose $p_n\to p$ and $p_n\ne p$.  If constants $\alpha>0$ and $\lambda>0$ exist such that
\[
 \lim_{n\to\infty}
 \frac{|p_{n+1}-p|}{|p_n-p|^\alpha}
 =\lambda,
\]
then the sequence converges to $p$ with order $\alpha$ and asymptotic error constant $\lambda$.
\end{definition}

Convergence is called linear when $\alpha=1$, superlinear when $1<\alpha<2$, and quadratic when $\alpha=2$.

\subsection{Aitken's $\Delta^2$ Process}

For a linearly convergent sequence $\{p_n\}$, Aitken's process forms
\[
 \widetilde p_n
 =p_n-\frac{(p_{n+1}-p_n)^2}
 {p_{n+2}-2p_{n+1}+p_n},
\]
provided the denominator is nonzero.  The transformed sequence often converges substantially faster.

\subsection{Steffensen's Method}

Steffensen's method applies Aitken acceleration directly to fixed-point iteration.  Starting with $p_0$, compute
\[
 p_1=g(p_0),
 \qquad
 p_2=g(p_1),
\]
and then
\[
 p=p_0-\frac{(p_1-p_0)^2}{p_2-2p_1+p_0}.
\]
Replace $p_0$ by $p$ and repeat.

\begin{example}
Use Steffensen's method to solve
\[
 x^3+4x^2-10=0,
 \qquad
 g(x)=\sqrt{\frac{10}{x+4}},
\]
starting from $p_0=1.5$.  Keep five decimal places.
\end{example}

\begin{solution}[4.6cm]
For the first cycle,
\[
 p_1=g(1.5)\approx1.34840,
 \qquad
 p_2=g(p_1)\approx1.36738.
\]
Aitken acceleration gives
\[
 p\approx1.36526.
\]
Repeating the process and rounding to five decimal places gives
\[
 p\approx1.36523.
\]
\end{solution}

\section{Zeros of Polynomials}

\begin{theorem}[Fundamental theorem of algebra]
Every polynomial of degree $n\ge1$ with real or complex coefficients has at least one complex root.  Counting multiplicity, it has exactly $n$ complex roots.
\end{theorem}

\subsection{Müller's Method}

Müller's method fits a quadratic polynomial through three points and uses one root of that quadratic as the next approximation.  It can naturally generate complex approximations even when all starting values are real.

Given $x_0,x_1,x_2$, write the local quadratic in the form
\[
 q(x)=a(x-x_2)^2+b(x-x_2)+c,
 \qquad c=f(x_2).
\]
The coefficients $a$ and $b$ satisfy
\[
\begin{pmatrix}
(x_0-x_2)^2&x_0-x_2\\
(x_1-x_2)^2&x_1-x_2
\end{pmatrix}
\begin{pmatrix}a\\b\end{pmatrix}
=
\begin{pmatrix}f(x_0)-c\\f(x_1)-c\end{pmatrix}.
\]
Then
\[
 x_3=x_2+\frac{-2c}{b\pm\sqrt{b^2-4ac}},
\]
where the sign is chosen to make the denominator as large as possible in magnitude.

\begin{example}
Use one Müller iteration for
\[
 f(x)=x^3-2x^2-5,
 \qquad x_0=1,
 \quad x_1=2,
 \quad x_2=3.
\]
\end{example}

\begin{solution}[5.0cm]
The function values are
\[
 f(1)=-6,
 \qquad f(2)=-5,
 \qquad f(3)=4.
\]
Thus $c=4$, and solving
\[
\begin{pmatrix}4&-2\\1&-1\end{pmatrix}
\begin{pmatrix}a\\b\end{pmatrix}
=
\begin{pmatrix}-10\\-9\end{pmatrix}
\]
gives $a=4$ and $b=13$.  Therefore
\[
 x_3=3+\frac{-8}{13+\sqrt{105}}
 \approx2.6559.
\]
\end{solution}

\section{Systems of Nonlinear Equations}

Let
\[
 \vect F(\vect X)=
 \begin{pmatrix}
 f_1(x_1,\ldots,x_n)\\
 \vdots\\
 f_n(x_1,\ldots,x_n)
 \end{pmatrix}.
\]
Newton's method for $\vect F(\vect X)=\vect0$ uses the Jacobian matrix
\[
 J(\vect X)=
 \left[\frac{\partial f_i}{\partial x_j}\right].
\]
At iteration $k$, solve
\[
 J(\vect X^{(k)})\vect H^{(k)}
 =-\vect F(\vect X^{(k)})
\]
and update
\[
 \vect X^{(k+1)}=\vect X^{(k)}+\vect H^{(k)}.
\]

\begin{example}
Apply one Newton iteration to
\[
\begin{aligned}
 x_1^2-x_2^2+2x_2&=0,\\
 2x_1+x_2^2-6&=0,
\end{aligned}
\]
starting from $(x_1,x_2)=(0.6,2.2)$.
\end{example}

\begin{solution}[5.5cm]
At $\vect X^{(0)}=(0.6,2.2)\T$,
\[
 \vect F(\vect X^{(0)})
 =\begin{pmatrix}-0.08\\0.04\end{pmatrix},
 \qquad
 J(\vect X^{(0)})
 =\begin{pmatrix}1.2&-2.4\\2&4.4\end{pmatrix}.
\]
Solve
\[
 \begin{pmatrix}1.2&-2.4\\2&4.4\end{pmatrix}
 \begin{pmatrix}h_1\\h_2\end{pmatrix}
 =\begin{pmatrix}0.08\\-0.04\end{pmatrix}.
\]
This gives
\[
 h_1\approx0.0254,
 \qquad
 h_2\approx-0.0206.
\]
Therefore
\[
 \vect X^{(1)}
 \approx\begin{pmatrix}0.6254\\2.1794\end{pmatrix}.
\]
\end{solution}
